% Source: https://artofproblemsolving.com/community/c3h2358851_easy_question
% Stable identity: aops-topic-2358851-q01
Number the places $0,1,\dots,23$ around the table. Rotating the table by $r$ places ($r=0,1,\dots,23$) moves every card $r$ places forward, so for each guest exactly one of the $24$ rotations brings his or her card in front of him or her.

\textbf{(a)} The rotation by $0$ places serves no guest, since nobody is correctly seated. So the $24$ guests are distributed among the $23$ rotations $r=1,\dots,23$, and by the pigeonhole principle some rotation serves at least $2$ guests: after that rotation, at least two guests sit in front of their own cards. $\blacksquare$

\textbf{(b)} \textbf{No.} Seat guest $i$ at place $i$ for $i=0,1,\dots,23$, and put the card of guest $i$ at place $1-i$, reduced modulo $24$ (so the cards of guests $0,1,2,\dots,23$ are at places $1,0,23,\dots,2$). Different guests get different places, so this is a valid placement of the cards.
\begin{itemize}
\item Nobody is correctly seated at the start: $1-i\equiv i\pmod{24}$ would mean $2i\equiv1\pmod{24}$, but $2i-1$ is odd and cannot be a multiple of $24$.
\item After a rotation by $r$ places, the card of guest $i$ is at place $1-i+r$, so guest $i$ is correctly seated exactly when $2i\equiv r+1\pmod{24}$. If $r$ is even, $r+1$ is odd and no guest is correctly seated. If $r$ is odd, the condition means $i\equiv\dfrac{r+1}{2}\pmod{12}$, which holds for exactly two of the numbers $i=0,1,\dots,23$.
\end{itemize}
So every rotation seats either $0$ or exactly $2$ guests correctly, and three guests can never be correctly seated at the same time. $\blacksquare$
